% GENERATED FILE. Edit canonical lessons, then run npm run generate:books.
% canonical-source-hash: fnv1a64:867766db07a945d4
% canonical-lessons: FA-C277-valedeyn, FA-C277-namzad, FA-C277-damad, FA-C277-famil, FA-C277-nojavan

\chapter{Parents, Fiancé, Bridegroom, Relatives, Teenager}
\label{ch:fa-parents-fiance-bridegroom-relatives-teenager}
\hlchaptermodality{277}

\begin{chapteropening}
\textbf{By the end of this chapter:} \emph{I can say parents, a fiancé, a bridegroom, relatives and a teenager.}

Complete the last lesson of chapter 277: I can say parents, a fiancé, a bridegroom, relatives and a teenager.
\end{chapteropening}

\section[\texorpdfstring{\fa{والدین} (vâledeyn)}{والدین (vâledeyn)}]{\fa{والدین} (vâledeyn) — parents}

\label{lesson:FA-C277-valedeyn}



\begin{quote}
Before the new one: say the Persian for a text message, then the Persian for an email.
\end{quote}

\subsection*{You'll want to know: \fa{والدین}}
\textbf{\fa{والدین}} — \emph{vâledeyn} — \textquotedblleft{}parents\textquotedblright{}.

\begin{cousinweb}
\emph{vâledeyn} is borrowed from Arabic \emph{wālidayn}, \textquotedblleft{}the two parents\textquotedblright{}. Its Arabic root is \emph{w-l-d} (Lane). It is the Arabic dual form, \textquotedblleft{}two parents\textquotedblright{}.
\end{cousinweb}

\begin{grammarlens}[title={What you've built}]
One more word for this chapter.
\end{grammarlens}

\subsection*{Your turn}
\begin{itemize}
\raggedright
  \item \emph{Say it:} \emph{vâledeyn}
  \item \emph{Say it:} \emph{vâledeyn}, once more
  \item \emph{Say it:} say \emph{vâledeyn}
  \item \emph{Recall:} say the Persian for a text message, then the Persian for an email, then say \emph{vâledeyn} again
\end{itemize}

\begin{tcolorbox}[breakable,colback=teal!4,colframe=teal!35!black,title={Before you move on}]
What does \fa{والدین} mean? (\textquotedblleft{}Parents\textquotedblright{}.) Say it once more.
\end{tcolorbox}
\section[\texorpdfstring{\fa{نامزد} (nâmzad)}{نامزد (nâmzad)}]{\fa{نامزد} (nâmzad) — a fiancé, a fiancée}

\label{lesson:FA-C277-namzad}



\begin{quote}
Before the new one: say the Persian for an email, then the Persian for parents.
\end{quote}

\subsection*{You'll want to know: \fa{نامزد}}
\textbf{\fa{نامزد}} — \emph{nâmzad} — \textquotedblleft{}a fiancé, a fiancée\textquotedblright{}.

\begin{cousinweb}
\emph{nâmzad} is \emph{nâm}, \textquotedblleft{}name\textquotedblright{} (inherited, EtymDB), with \emph{zad}, \textquotedblleft{}struck\textquotedblright{} (from \emph{zadan}): one whose name has been set down for another. Both are native Persian.
\end{cousinweb}

\begin{grammarlens}[title={What you've built}]
One more word for this chapter.
\end{grammarlens}

\subsection*{Your turn}
\begin{itemize}
\raggedright
  \item \emph{Say it:} \emph{nâmzad}
  \item \emph{Say it:} \emph{nâmzad}, once more
  \item \emph{Say it:} say \emph{nâmzad}
  \item \emph{Recall:} say the Persian for an email, then the Persian for parents, then say \emph{nâmzad} again
\end{itemize}

\begin{tcolorbox}[breakable,colback=teal!4,colframe=teal!35!black,title={Before you move on}]
What does \fa{نامزد} mean? (\textquotedblleft{}A fiancé, a fiancée\textquotedblright{}.) Say it once more.
\end{tcolorbox}
\section[\texorpdfstring{\fa{داماد} (dâmâd)}{داماد (dâmâd)}]{\fa{داماد} (dâmâd) — a bridegroom; a son-in-law}

\label{lesson:FA-C277-damad}



\begin{quote}
Before the new one: say the Persian for parents, then the Persian for a fiancé.
\end{quote}

\subsection*{You'll want to know: \fa{داماد}}
\textbf{\fa{داماد}} — \emph{dâmâd} — \textquotedblleft{}a bridegroom; a son-in-law\textquotedblright{}.

\begin{cousinweb}
\emph{dâmâd} is a native Persian word, inherited rather than borrowed: it continues Middle Persian \emph{dāmād}, \textquotedblleft{}bridegroom, son-in-law\textquotedblright{} (EtymDB). EtymDB takes it back to the Indo-European root \emph{*ǵem-}, \textquotedblleft{}to marry\textquotedblright{}, the root of Greek \emph{gambros}, \textquotedblleft{}son-in-law\textquotedblright{}.
\end{cousinweb}

\begin{grammarlens}[title={What you've built}]
One more word for this chapter.
\end{grammarlens}

\subsection*{Your turn}
\begin{itemize}
\raggedright
  \item \emph{Say it:} \emph{dâmâd}
  \item \emph{Say it:} \emph{dâmâd}, once more
  \item \emph{Say it:} say \emph{dâmâd}
  \item \emph{Recall:} say the Persian for parents, then the Persian for a fiancé, then say \emph{dâmâd} again
\end{itemize}

\begin{tcolorbox}[breakable,colback=teal!4,colframe=teal!35!black,title={Before you move on}]
What does \fa{داماد} mean? (\textquotedblleft{}A bridegroom; a son-in-law\textquotedblright{}.) Say it once more.
\end{tcolorbox}
\section[\texorpdfstring{\fa{فامیل} (fâmil)}{فامیل (fâmil)}]{\fa{فامیل} (fâmil) — relatives, extended family}

\label{lesson:FA-C277-famil}



\begin{quote}
Before the new one: say the Persian for a fiancé, then the Persian for a bridegroom.
\end{quote}

\subsection*{You'll want to know: \fa{فامیل}}
\textbf{\fa{فامیل}} — \emph{fâmil} — \textquotedblleft{}relatives, extended family\textquotedblright{}.

\begin{cousinweb}
\emph{fâmil} is borrowed from French \emph{famille}, \textquotedblleft{}family\textquotedblright{} (EtymDB).
\end{cousinweb}

\begin{grammarlens}[title={What you've built}]
One more word for this chapter.
\end{grammarlens}

\subsection*{Your turn}
\begin{itemize}
\raggedright
  \item \emph{Say it:} \emph{fâmil}
  \item \emph{Say it:} \emph{fâmil}, once more
  \item \emph{Say it:} say \emph{fâmil}
  \item \emph{Recall:} say the Persian for a fiancé, then the Persian for a bridegroom, then say \emph{fâmil} again
\end{itemize}

\begin{tcolorbox}[breakable,colback=teal!4,colframe=teal!35!black,title={Before you move on}]
What does \fa{فامیل} mean? (\textquotedblleft{}Relatives, extended family\textquotedblright{}.) Say it once more.
\end{tcolorbox}
\section[\texorpdfstring{\fa{نوجوان} (nojavân)}{نوجوان (nojavân)}]{\fa{نوجوان} (nojavân) — a teenager}

\label{lesson:FA-C277-nojavan}



\begin{quote}
Before the new one: say the Persian for a bridegroom, then the Persian for relatives.
\end{quote}

\subsection*{You'll want to know: \fa{نوجوان}}
\textbf{\fa{نوجوان}} — \emph{nojavân} — \textquotedblleft{}a teenager\textquotedblright{}.

\begin{cousinweb}
\emph{nojavân} is \emph{no}, \textquotedblleft{}new\textquotedblright{}, with \emph{javân}, \textquotedblleft{}young\textquotedblright{}: newly young. Both are inherited; EtymDB takes \emph{javân} back to the Indo-European word behind Latin \emph{iuvenis}.
\end{cousinweb}

\begin{grammarlens}[title={What you've built}]
One more word for this chapter.
\end{grammarlens}

\subsection*{Your turn}
\begin{itemize}
\raggedright
  \item \emph{Say it:} \emph{nojavân}
  \item \emph{Say it:} \emph{nojavân}, once more
  \item \emph{Say it:} say \emph{nojavân}
  \item \emph{Recall:} say the Persian for a bridegroom, then the Persian for relatives, then say \emph{nojavân} again
\end{itemize}

\begin{tcolorbox}[breakable,colback=teal!4,colframe=teal!35!black,title={Before you move on}]
What does \fa{نوجوان} mean? (\textquotedblleft{}A teenager\textquotedblright{}.) Say it once more.
\end{tcolorbox}
